\section{Structural proof obligations}
Core proof/test obligations include generation non-authority, mechanism non-upgrade, typed relation non-escalation, non-forced gluing, negative-history monotonicity, conservative evidence-lineage saturation, missing-evidence honesty, residual reopening, meta-evaluator separation, measurement non-escalation, explicit uncertainty assumptions and formal-closure honesty. These are structural constraints on the method; they do not prove empirical scientific performance.

\begingroup
\small
\setlength{\emergencystretch}{3em}
\subsection{Additional obligations exposed by the saturation pass}

The deeper comparison adds several obligations that were implicit in earlier versions.

\paragraph{Geometry obligation.}
Any future claim that a \RAKL{} object is an order-theoretic lattice must name the scoped order and demonstrate the required meet/join or closure-system properties. Any sheaf claim must name the local objects, restriction maps and gluing condition. Navigation-graph adjacency and J-space vector geometry do not satisfy those obligations by analogy.

\paragraph{Cognition/evidence separation obligation.}
A workspace intervention or interpretability observation can establish computational influence only. Promoting it to scientific evidence requires an explicit bridge to the external scientific object and the normal verification path.

\paragraph{Manuscript-saturation obligation.}
A paper-level saturation certificate must report its source universe, freshness cutoff, route coverage, review lenses, last material semantic additions and remaining open classifications. Same-context flatness must never be described as independent peer review.

\paragraph{Novelty-reopen obligation.}
Discovery of prior work that semantically subsumes an integrated \RAKL{} claim must narrow the novelty claim even if the terminology differs. The prior search transcript remains negative method history and the affected saturation certificate is reopened.

\paragraph{Evaluation obligation.}
Scientific-agent superiority cannot be inferred from software-test count, toy-world correctness, context compression or same-context self-repair. It requires the preregistered matched and fresh evaluations under the published resource and authority criteria.

\subsection{Non-equivalence obligations}

Several negative claims in the paper are structural and should be read as non-equivalence obligations rather than empirical performance claims.

\paragraph{Graph is not authority.}
\RAKL{} defines no generic authority-preserving map from note-link degree, citation count or retrieval centrality to scientific authority. Any proposed map must separately show that it preserves every blocking coordinate of the scoped authority poset. A visualization may correlate with relevance while remaining non-authoritative.

\paragraph{Access is not evidence.}
The framework defines no generic authority-preserving map from J-space occupancy or workspace salience to external scientific truth. Such a map requires an additional empirical bridge specific to the scientific object. The default relation is cognitive provenance only.

\paragraph{Compatibility is not a lattice.}
Pairwise or higher-order co-admissibility does not supply an order, meet or join. An order-theoretic lattice claim requires separately verified structure; FCA illustrates one legitimate route \cite{ganterwille2024}.

\paragraph{Global existence is not identification.}
The existence of a global section or coherent global model does not imply uniqueness, and uniqueness of a mathematical representation does not imply evidence sufficient for scientific authority. These predicates live on different certificates.

\paragraph{Traceability is not independence.}
Two results with complete provenance can still share one evidence root. Independent-corroboration credit therefore requires evidence-lineage analysis in addition to ordinary provenance.

These obligations are useful because they define simple reviewer attacks. A counterexample that violates one non-equivalence under the claimed scope is enough to reopen the corresponding part of the framework. They are also implementation tests: any code path that raises authority because an item was frequently retrieved, highly linked or present in a workspace violates the architecture even if downstream answers improve.
\endgroup
